\documentclass{article}

\usepackage{amsmath, amssymb, amsthm}
\usepackage[a4paper, margin=1in]{geometry}

\title{Fermat Primality Test}
\author{sid}
\date{\today}

\newtheorem{theorem}{Theorem}
\newtheorem{lemma}{Lemma}
\newtheorem{example}{Example}

\newcommand{\Z}{\mathbb{Z}}
\newcommand{\bld}[1]{\textbf{#1}}
\newcommand{\itl}[1]{\textit{#1}}
\newcommand{\N}{\mathbb{N}}
\newcommand{\Q}{\mathbb{Q}}
\newcommand{\R}{\mathbb{R}}
\newcommand{\eps}{\varepsilon}
\newcommand{\poly}{\mathrm{poly}}
\newcommand{\polylog}{\mathrm{polylog}}
\newcommand{\bigO}{\mathcal{O}}

\begin{document}

\maketitle
The simplest probabilistic primality test is the Fermat primality test (actually a compositeness test). It works as follows:

\vspace{1em}
If $a^{n-1} \not\equiv 1 \pmod{n}$, then $n$ is composite.
where, $a$ is any co-prime of $n$

If $a^{n-1}$(modulo $n$) is 1 but $n$ is not a prime, then $n$ is called a \itl{pseudoprime to base $a$}
\begin{example}
  Let $n = 341$ and $a = 2$. Then
  \[
  2^{340} \equiv 1 \pmod{341},
  \]
  even though $341 = 11 \cdot 31$ is composite.
  In fact, 341 is the smallest pseudoprime to the base of 2.
\end{example}

\begin{theorem}[Fermat's Little Theorem]
Let $p$ be a prime and let $a \in \mathbb{Z}$ with $p \nmid a$. Then
\[
a^{p-1} \equiv 1 \pmod{p}.
\]
\end{theorem}

\begin{proof}
Consider the set of nonzero residues modulo $p$:
\[
S = \{1,2,3,\dots,p-1\}
\]
Now consider the set obtained by multiplying each element of $S$ by $a$:
\[
aS = \{a,2a,3a,\dots,(p-1)a\}
\]

\bld{Claim:} $aS$ is just a rearrangement of $S$

\itl{Step 1: all elements of aS are unique}
For contracdiction, let us assume there exist two distinct $i$ and $j$ in $S$ such that
\[
\begin{aligned}
ai &\equiv aj \pmod{p} \\
\implies ai - aj &\equiv 0 \pmod{p} \\
\implies a(i - j) &\equiv 0 \pmod{p}.
\end{aligned}
\]

Since $p \nmid a$, we have $\gcd(a,p)=1$, so $a$ is invertible modulo $p$. Hence,
\[
i - j \equiv 0 \pmod{p}.
\]

Since $1 \le i,j \le p-1$, this implies $-p<i-j<p$. Then this is only possible when $i=j$, i.e., $i$ and $j$ are not distinct.

\itl{Step 2: all elements of aS are non-zero}
This is very simple to prove as neither $a$ or $i\in S$ are divisible by $p$ so $ai$ cannot be divisible by $p$

\bld{Claim proved.} as we have $p-1$ unique elements all less than $p$ and none of them are 0

Multiplying each element in $aS$,

\[
\begin{aligned}
  a\cdot2a\cdot3a\dots\cdot{(p-1)a} &\equiv 1\cdot2\cdot3\dots\cdot{(p-1)} \pmod p \\
  a^{p-1}\cdot(p-1)! &\equiv  (p-1)! \pmod p \\
  a^{p-1}&\equiv  1\pmod p
\end{aligned}
\]
\end{proof}

\section*{Time Complexity}
The dominant operation is computing $a^{n-1} \bmod n$. This is done using
\emph{fast modular exponentiation} (repeated squaring).

The exponent $n-1$ has $O(\log n)$ bits, so modular exponentiation requires
$O(\log n)$ modular multiplications.

Each multiplication involves numbers of size at most $n$, i.e., with $O(\log n)$ bits.
Using schoolbook multiplication, this costs
\[
O((\log n)^2).
\]

Thus, one iteration of the Fermat test runs in
\[
O((\log n)^3).
\]

To improve reliability, the test is repeated with different choices of $a$.
If we perform $k$ independent trials, the total running time is
\[
O(k (\log n)^3).
\]

Using faster multiplication algorithms (e.g., Karatsuba or FFT-based methods),the complexity can be improved to
\[
O\big(k \log^2 n \log \log n\big).
\]

The Fermat primality test runs in polynomial time in the input size $\log n$, making it computationally efficient.

\section*{Flaws}

Although the Fermat primality test is computationally efficient, it suffers from
fundamental limitations that make it unreliable in practice.

\subsection*{Pseudoprimes}

There are only 21853 pseudoprimes base 2 that are less than $2.5\times10^{10}$. This means that for $n$ upto $2.5\times10^{10}$, Fermat's test fails only $8.7\times{10^{-7}}\%$ of times, 


\subsection*{Carmichael Numbers}

A more serious issue arises from a special class of composite numbers.

\bld{Definition} A composite integer $n$ is called a \emph{Carmichael number} if

\[
a^{n-1} \equiv 1 \pmod{n}
\quad \text{for all integers } a \text{ with } \gcd(a,n)=1.
\]
In other words, Carmichael numbers satisfy Fermat's congruence for \emph{every} valid base $a$, not just some.
\begin{theorem}[Korselt's Criterion]
A composite integer $n$ is a Carmichael number if and only if:
\begin{itemize}
    \item $n$ is square-free (No perfect square greater than 1 divides $n$), and
    \item for every prime divisor $p \mid n$, we have
    \[
    p - 1 \mid n - 1.
    \]
\end{itemize}
\end{theorem}

This provides a constructive way to identify Carmichael numbers.

\begin{example}
  Consider
  \[
  n = 561 = 3 \cdot 11 \cdot 17.
  \]

  We verify Korselt's criterion:
  \[
  \begin{aligned}
  3 - 1 &= 2 \mid 560, \\
  11 - 1 &= 10 \mid 560, \\
  17 - 1 &= 16 \mid 560.
  \end{aligned}
  \]

  Since $561$ is square-free and satisfies the divisibility conditions, it is a
  Carmichael number. Consequently,
  \[
  a^{560} \equiv 1 \pmod{561}
  \]
  for all $a$ with $\gcd(a,561)=1$.
\end{example}

\section*{Conclusion}

The Fermat primality test is fast and simple, but unreliable due to the existence of pseudoprimes and Carmichael numbers. While useful as a preliminary filter, it cannot serve as a definitive test for primality.

\end{document}